\section{Minors, cofactors, and Laplace expansion}

Sarrus' rule gives a quick method for \(3\times3\) matrices, but it does not generalize to larger matrices. If we want a method that works for every square matrix, we need a different idea. The key idea is reduction: we compute a large determinant by reducing it to smaller determinants.

This is a recurring mathematical pattern. When an object is too large to understand at once, we try to break it into smaller objects of the same kind. For determinants, the smaller objects are called minors.

\subsection{Minors}

Let \(A=(a_{ij})\) be an \(n\times n\) matrix. Choose an entry \(a_{ij}\). Remove row \(i\) and column \(j\). The determinant of the remaining \((n-1)\times(n-1)\) matrix is called the minor of \(a_{ij}\), and is denoted by \(M_{ij}\).

\begin{examplebox}
Let
\[
A=\begin{pmatrix}
2&1&3\\
0&-1&4\\
5&2&1
\end{pmatrix}.
\]
To find \(M_{12}\), we remove row \(1\) and column \(2\). The remaining matrix is
\[
\begin{pmatrix}
0&4\\
5&1
\end{pmatrix}.
\]
Therefore
\[
M_{12}=\det\begin{pmatrix}0&4\\5&1\end{pmatrix}=0\cdot1-4\cdot5=-20.
\]
\end{examplebox}

A minor is not an entry of the original matrix. It is a determinant obtained after deleting a row and a column. This distinction matters: the entry \(a_{ij}\) is a number already present in the matrix; the minor \(M_{ij}\) is computed from what remains after removing the row and column of that entry.

\subsection{Cofactors}

The minor still needs a sign. The signs follow a checkerboard pattern:
\[
\begin{pmatrix}
+&-&+&- & \cdots\\
-&+&-&+ & \cdots\\
+&-&+&- & \cdots\\
-&+&-&+ & \cdots\\
\vdots&\vdots&\vdots&\vdots&\ddots
\end{pmatrix}.
\]
This pattern is encoded by the factor \((-1)^{i+j}\).

\begin{definitionbox}
The cofactor of the entry \(a_{ij}\) is
\[
C_{ij}=(-1)^{i+j}M_{ij}.
\]
\end{definitionbox}

\begin{examplebox}
Using the matrix from the previous example, we found
\[
M_{12}=-20.
\]
Since \((-1)^{1+2}=(-1)^3=-1\), the cofactor is
\[
C_{12}=(-1)M_{12}=20.
\]
For the entry in position \((1,1)\), the sign is positive. Removing row \(1\) and column \(1\), we get
\[
M_{11}=\det\begin{pmatrix}-1&4\\2&1\end{pmatrix}=(-1)\cdot1-4\cdot2=-9.
\]
Thus
\[
C_{11}=M_{11}=-9.
\]
\end{examplebox}

The alternating signs are not decoration. Without them, the determinant would not have the correct behaviour under row exchange, and it would not detect orientation properly.

\subsection{Laplace expansion along a row}

Laplace expansion expresses a determinant as a sum of entries multiplied by their cofactors. If we expand along row \(i\), then
\[
\det A=a_{i1}C_{i1}+a_{i2}C_{i2}+\cdots+a_{in}C_{in}.
\]
Equivalently,
\[
\det A=\sum_{j=1}^{n}a_{ij}C_{ij}.
\]

\begin{examplebox}
Let
\[
A=\begin{pmatrix}
2&1&3\\
0&-1&4\\
5&2&1
\end{pmatrix}.
\]
Expand along the first row. We need the cofactors \(C_{11}\), \(C_{12}\), and \(C_{13}\).

First,
\[
C_{11}=+\det\begin{pmatrix}-1&4\\2&1\end{pmatrix}
=(-1)\cdot1-4\cdot2=-9.
\]
Second,
\[
C_{12}=-\det\begin{pmatrix}0&4\\5&1\end{pmatrix}
=-(0\cdot1-4\cdot5)=20.
\]
Third,
\[
C_{13}=+\det\begin{pmatrix}0&-1\\5&2\end{pmatrix}
=0\cdot2-(-1)\cdot5=5.
\]
Therefore
\[
\det A=2C_{11}+1C_{12}+3C_{13}
=2(-9)+1(20)+3(5).
\]
Hence
\[
\det A=-18+20+15=17.
\]
\end{examplebox}

Notice how the computation works. A \(3\times3\) determinant was reduced to three \(2\times2\) determinants. For a \(4\times4\) determinant, Laplace expansion reduces the problem to \(3\times3\) determinants, and so on. The method is recursive.

\subsection{Laplace expansion along a column}

We may also expand along a column. If we expand along column \(j\), then
\[
\det A=a_{1j}C_{1j}+a_{2j}C_{2j}+\cdots+a_{nj}C_{nj}.
\]
Equivalently,
\[
\det A=\sum_{i=1}^{n}a_{ij}C_{ij}.
\]

The freedom to choose a row or a column is important. A wise choice can make a determinant much easier to compute.

\begin{examplebox}
Consider
\[
B=\begin{pmatrix}
4&0&1&2\\
0&3&0&-1\\
2&0&5&0\\
1&0&0&6
\end{pmatrix}.
\]
The second column contains three zeros:
\[
\begin{pmatrix}0\\3\\0\\0\end{pmatrix}.
\]
So we expand along the second column. Only one term survives:
\[
\det B=3C_{22}.
\]
Since \((-1)^{2+2}=1\),
\[
C_{22}=\det\begin{pmatrix}
4&1&2\\
2&5&0\\
1&0&6
\end{pmatrix}.
\]
Now compute the remaining \(3\times3\) determinant. Using Sarrus' rule,
\[
\det\begin{pmatrix}
4&1&2\\
2&5&0\\
1&0&6
\end{pmatrix}
=(4\cdot5\cdot6+1\cdot0\cdot1+2\cdot2\cdot0)
-(2\cdot5\cdot1+1\cdot2\cdot6+4\cdot0\cdot0).
\]
Thus
\[
C_{22}=120-(10+12)=98.
\]
Therefore
\[
\det B=3\cdot98=294.
\]
\end{examplebox}

This example shows why Laplace expansion is not merely a theoretical formula. It gives a strategy. Before computing, inspect the matrix. A row or column with many zeros is usually the best place to expand.

\subsection{A method, not a ritual}

Laplace expansion can be used along any row or any column. The result is always the same determinant. But the amount of work may be very different. The mathematical question is not only ``Can I compute it?'' but also ``How can I compute it intelligently?''

\begin{summarybox}
When using Laplace expansion, follow this order.
\begin{itemize}
\item Inspect the matrix before calculating.
\item Look for a row or column with many zeros.
\item Write the sign pattern carefully.
\item Compute the minors as smaller determinants.
\item Keep the structure visible; do not compress too many steps at once.
\end{itemize}
\end{summarybox}
